\scanpage{54}
\mspage{53}

Mais dans le cas polyédral, il faut un axiome suppl.

\[
\begin{aligned}
&\text{\uline{Prat pol}}\quad \forall X\in\unsure{M|\mathcal M},\quad \widetilde X\ \text{est polyédral,}\quad \text{i.e. si } \deleted{\unsure{X}}Y,Z\in M\ \text{et}\\
&\left\lvert\begin{aligned}
&\text{si }\{X,Y\}\ \text{est majoré et minoré dans }M,\ \text{alors }\mathrm{Inf}^{\leq}(\deleted{\unsure{X}}Y,\deleted{\unsure{Y}}Z)\ \text{existe}\\
&(\text{i.e. }\widetilde X\cap\widetilde Y\ \text{a un plus grand élément}).
\end{aligned}\right.
\end{aligned}
\]

\uline{NB} Alors $\mathrm{Sup}^{\leq}(\deleted{\unsure{X}}Y,\deleted{\unsure{Y}}Z)$ existe aussi \uline{dans} $\widetilde X$ (mais prob. \unsure{bien|sûr} \uline{pas} dans $M$\ldots)

\uline{Remarque} Soit $(M,\leq)$ un ens. ordonné, et \unsure{prenons|posons} $(\ll)\Leftrightarrow(\leq)$, $(x\parallel_{\unsure{\mathcal L|\mathscr L}}y)\Leftrightarrow(x\ne y)$. Alors tous les axiomes sont satisfaits, sauf le seul axiome de divisibilité $\mathrm{Prat}\,\mathcal L\,2'$. Je pense p.ex. au cas où $M={}\text{ensemble des simplexes}$ d'une maquette simpliciale, $(\unsure{\mathscr S|S},M\subset\unsure{P|\mathcal P}_{f}(\mathscr S))$, où les figures sont les sous-maquettes simpliciales, \unsure{correspondants|correspondant} aux sous-ensembles fermés de $(M,\leq(\overset{\mathrm{déf}}{=}\subset))$.

En termes de l'axiomatique \textcircled{1} ou \textcircled{2}, on prendrait $X\ll Y\deleted{\unsure{=} }\Longleftrightarrow X\leq Y$, et $X\mathrel{>\mkern-13mu<}Y$ tjrs satisfaits. \deleted{Alors $X\overset{\circ}{\ll}Y$ signifie $X=Y$.}

Appelons ``trivial'' un \inserted{\unsure{magasin|magasin,}} pré-atelier tel que

\[
X\ll Y\Longleftrightarrow X\leq Y\qquad X\mathrel{>\mkern-13mu<}Y\ \text{tjrs satisfaits.}
\]

de sorte que

\[
X\overset{\circ}{\ll}Y\Longleftrightarrow X=Y
\]

Alors tous les axiomes \inserted{Prat 1 - Prat 6} sont automatiquement satisfaits -- on peut donc partir d'un \deleted{\unsure{At|tel}} ens ordonné $(M,\leq)$ quelconque. Passant à $\unsure{\mathcal F|\mathscr F|F}$, on trouve \deleted{\unsure{c'est-à-dire|la totalité}}

\[
\begin{aligned}
&\mathcal F=P\!f(M,\leq)\qquad \text{ensemble de \uline{toutes}}\\
&\hspace{15em}\text{les parties fermées}
\end{aligned}
\]

% slanted
\margin{NB On trouve même un atelier \uline{séparé}, et \deleted{\unsure{disjoint}}\inserted{à multistrates près} \deleted{\illegible{1}} \unsure{ils|s'} sont (sauf erreur) \unsure{contrts|contr.} \unsure{//éditiles|//éditibles|//édités} etc\ldots}

% slanted
\margin{Pourtant, on peut ``faire de la topologie'' \unsure{cl$^{\mathrm{que}}$|classique|clque} sur une maquette simpliciale !}

dans $\mathcal F$, $F\ll G\Longleftrightarrow F\leq G$, et $F\mathrel{>\mkern-13mu<}G$ tjrs satisfait $F\parallel G\Longleftrightarrow F\cap G=\varnothing$. \deleted{\unsure{Axiome} C \unsure{est}} Les conditions $\mathrm{At}\,D$, $\mathrm{At}\,C(\unsure{\mathrm{spac}|\mathrm{spa}})$, $\mathrm{At}\,\mathrm{Fil}$, $\mathrm{At}\,\mathrm{supp}$ sont satisfaites\unsure{;|.} Pour ce \deleted{$\unsure{\mathrm{At}\,L\,1}$} qui est des conditions sur les lieux, $\mathrm{At}\,L\,2$ équivaut :

Tout \deleted{$\forall x\in M$} $X\in M$ est minoré par un élément minimal -- \unsure{c'est|n'est} pas \unsure{prédonnant|prédonné|précédent|prédécidé} ! L'axiome $\mathrm{At}\,L\,3$ tjrs satisfait (car deux lieux distincts sont disjoints, étant compatibles).

Donc le seul axiome pas satisfait est $\mathrm{At}\,L\,1$, divisibilité

